\chapter{Phương pháp đánh giá}
\label{chap:phuong-phap-danh-gia}

Trong chương này, báo cáo sẽ trình bày chi tiết về thiết kế thực nghiệm nhằm kiểm chứng sự hiệu quả của phương pháp đã đề xuất. Nội dung của chương sẽ bao gồm giới thiệu về bộ dữ liệu được sử dụng trong thực nghiệm, quy trình tạo dữ liệu ảnh nhiễu tổng hợp cho bài toán OCR, cấu hình chi tiết của hệ thống và mô tả về các độ đo được sử dụng.

\section{Bộ dữ liệu}

\textbf{RACE (Reading Comprehension Dataset From Examinations)} \cite{lai2017race} là một bộ dữ liệu hỏi đáp có quy mô lớn trong lĩnh vực xử lý ngôn ngữ tự nhiên. Khác với các bộ dữ liệu hỏi đáp khác như \textbf{SQuAD} \cite{rajpurkar2016squad} được trích xuất chủ yếu từ Wikipedia, RACE được xây dựng trực tiếp từ bộ đề thi của các kỳ thi tiếng Anh thực tế dành cho học sinh trung học cơ sở và trung học phổ thông tại Trung Quốc.

Nghiên cứu sử dụng RACE để làm bộ dữ liệu đánh giá bởi hai lý do chính. Thứ nhất, các văn bản trong RACE bao trùm nhiều lĩnh vực trong đời sống như giáo dục, tin tức, văn hóa, chính trị và khoa học. Nhờ đó, văn phong và đặc điểm dữ liệu của RACE rất giống với dữ liệu thực tế. Theo thống kê, các văn bản trong bộ dữ liệu này có độ dài từ 300 đến 500 từ. Số lượng này đủ lớn để chứa được lượng thông tin phục vụ cho một bài kiểm tra như các chuỗi sự kiện phức tạp, luận điểm trong khi vẫn đảm bảo không gây tràn cửa sổ ngữ cảnh (context window) của các sLLM.

Trong phạm vi thực nghiệm, nghiên cứu lấy mẫu ngẫu nhiên 100 đoạn văn có độ dài lớn nhất từ tập con RACE-High làm dữ liệu đầu vào cho toàn bộ các thí nghiệm sinh câu hỏi. Bảng~\ref{tab:race_stats} trình bày thống kê tổng quan về bộ dữ liệu RACE và mẫu được sử dụng.

\section{Quy trình tổng hợp bộ dữ liệu OCR}

Bên cạnh bài toán sinh câu hỏi trắc nghiệm từ văn bản, nghiên cứu cần xây dựng một bộ dữ liệu nhằm đánh giá chất lượng của các VLM trong tác vụ OCR. Một thách thức lớn được đặt ra là các bộ dữ liệu OCR hiện tại không chuyên biệt cho các văn bản dài. Bên cạnh đó, việc thu thập và gán nhãn thủ công các ảnh chụp văn bản học liệu vô cùng tốn kém về chi phí và thời gian. Để giải quyết vấn đề, nghiên cứu triển khai một quy trình sinh dữ liệu tổng hợp dành riêng cho các ảnh chụp văn bản dài trong lĩnh vực giáo dục. Cụ thể, quy trình này sẽ sử dụng các văn bản từ bộ RACE để kết xuất thành các hình ảnh tương tự ảnh chụp màn hình. Sau đó các ảnh chụp này sẽ được thêm các nhiễu bằng cách mô phỏng tác nhân từ môi trường thực tế. Quy trình hai giai đoạn này được biểu diễn trong Hình~\ref{fig:augmentation-pipeline}.

\figure{}

\bold{Giai đoạn sinh ảnh sạch.} Hệ thống tạo ra các hình ảnh kỹ thuật số sắc nét từ văn bản đầu vào bằng cách sử dụng trình duyệt không giao diện (headless browser) từ thư viện Pyppeteer\cite{} để tận dụng sự linh hoạt của HTML và CSS. Quy trình này gồm 3 bước cốt lõi: Lựa chọn ngẫu nhiên một bộ cài đặt gồm kích thước ảnh, màu sắc, phông chữ (kết hợp cả phông chữ in chuẩn và viết tay từ Google Fonts\cite); tự động phân bổ mật độ hiển thị của văn bản như số cột, văn lề, khoảng cách giữa các dòng; và cuối cùng là nhúng mã Javascript để tự động co giãn cỡ chữ để toàn bộ văn bản nằm trọn trong khung hình.

\bold{Giai đoạn mô phỏng ảnh nhiễu.} Nghiên cứu sử dụng thư viện Augraphy\cite{} để triển khai quy trình tăng cường dữ liệu trên các hình ảnh sạch. Quá trình này tích hợp đa dạng các lớp biến đổi phức tạp bao gồm nhiễu do vật liệu giấy và thiết bị in ấn (vết mực nhòe, đứt nét, giấy ố vàng, chữ in chìm mặt sau), các nhiễu quang học phổ biến trong thực tế (phơi sáng, bóng đổ, lóa sáng) cũng như các nhiễu cơ học (giấy nhăn, gạch chân thủ công, lỗi nén ảnh). Bằng cách mô phỏng rất nhiều biến thể của nhiễu, nghiên cứu có thể đánh giá được khả năng trích xuất văn bản của các VLM trong các điều kiện thực tế khác nhau. Hình~\ref{fig:augmentation-samples} trình bày một số ảnh được làm nhiễu điển hình được sử dụng thực nghiệm.

\section{Thiết lập thực nghiệm}

Nghiên cứu đánh giá tổng cộng 12 mô hình được chia thành ba nhóm: nhóm các sLLM mã nguồn mở dùng để sinh câu hỏi trắc nghiệm, nhóm VLM quy mô nhỏ dùng cho tác vụ OCR và nhóm các LLM thương mại dành cho cả hai tác vụ.

Bảng~\ref{tab:models_qg} tóm tắt các mô hình sinh câu hỏi được sử dụng trong thực nghiệm. Các mô hình sLLM có kích thước tham số từ 8 đến 14 tỷ, phù hợp với mục tiêu triển khai trên phần cứng phổ thông. Các mô hình thương mại được đưa vào làm mô hình cơ sở để đánh giá mức độ cạnh tranh của giải pháp MAS kết hợp sLLM.

\table{gộp thành 1 bảng duy nhất của 12 mô hình}

Toàn bộ mã nguồn được viết bằng ngôn ngữ Python phiên bản 3.12. Thư viện {LangGraph} được sử dụng để cài đặt và triển khai kiến trúc đa tác tử một cách hiệu quả. Bên cạnh đó, để thuận tiện cho quá trình chạy thực nghiệm và chuyển đổi giữa các mô hình, nghiên cứu sử dụng dịch vụ API từ OpenRouter và NovitaAI để truy cập đến các mô hình sLLM, VLM và LLM thương mại.

Các siêu tham số chính được thiết lập như sau:
\begin{itemize}
    \item Nhiệt độ (temperature): 0,0 cho các tác tử chính (tác tử \textit{Sinh câu hỏi}, tác tử \textit{Phân loại}, tác tử \textit{Hiệu chỉnh}) nhằm đảm bảo tính nhất quán và tái lập được của kết quả; 0,3 cho tác tử \textit{Học sinh} nhằm tạo sự đa dạng trong cách trả lời.
    \item Số token tối đa (max\_tokens): 8.192 token cho mỗi lần gọi mô hình.
    \item Số vòng lặp tối đa (max\_loops): Giới hạn ở 8 vòng cho mỗi chu trình phản hồi giữa các tác tử.
\end{itemize}

Về quy mô của thực nghiệm, hệ thống được yêu cầu sinh $5$ câu hỏi cho mỗi cấp độ nhận thức, tổng cộng là $15$ câu hỏi cho mỗi đoạn văn đầu vào. Với số lượng là 100 đoạn văn mẫu thu thập ngẫu nhiên từ bộ RACE, mỗi mô hình cần sinh ra $1.500$ câu hỏi.

Để đánh giá hiệu quả của kiến trúc đa tác tử được đề xuất trong bài toán sinh câu hỏi trắc nghiệm, nghiên cứu tiến hành thiết lập hai kịch bản thử nghiệm trên mỗi mô hình sLLM: phương pháp cơ sở, yêu cầu mô hình tạo toàn bộ bài kiểm tra trong một lần gọi duy nhất (single prompt), và phương pháp tích hợp đa tác tử. Riêng đối với các LLM thương mại, nhằm đảm bảo tính công bằng về mặt tài nguyên đánh giá, nghiên cứu chỉ áp dụng cơ chế gọi mô hình một lần. Quyết định này được đưa ra dựa trên thực tế rằng chi phí vận hành cho một lượt truy vấn đối với các LLM thương mại thường tương đương, hoặc thậm chí vượt trội hơn so với tổng chi phí triển khai hệ thống sLLM đa tác tử. 

Đối với quy trình LLM-as-a-Judge, nghiên cứu sử dụng Gemini 3 Flash làm mô hình giám khảo cho tất cả thí nghiệm. Lý do nghiên cứu chọn mô hình này dựa vào ba tiêu chí chính: chất lượng suy luận cao cho các tác vụ đánh giá, tốc độ xử lý nhanh và chi phí hợp lý cho việc đánh giá quy mô lớn.

\section{Phương pháp và Tiêu chí Đánh giá}

Mục này sẽ trình bày chi tiết các độ đo được sử dụng trong nghiên cứu. Hệ thống triển khai hai khâu đánh giá độc lập cho hai bài toán tương ứng: đánh giá chất lượng câu hỏi trắc nghiệm và đánh giá độ chính xác của tác vụ OCR.

\subsection{Đánh giá chất lượng câu hỏi trắc nghiệm}

Đối với bài toán sinh câu hỏi trắc nghiệm, các thang đo NLP truyền thống như BLEU hay ROUGE tỏ ra kém hiệu quả và thiếu tin cậy do điểm yếu cốt lõi là chỉ so sánh sự trùng lặp bề mặt mà thiếu định sự thấu hiểu về ngôn ngữ và ngữ cảnh. Bằng cách áp dụng kỹ thuật LLM-as-a-Judge, hệ thống có thể đánh giá một cách chính xác và khác quan các tiêu chí sư phạm. Cụ thể, nghiên cứu sử dụng ba tiêu chí sư phạm sau:

\textbf{Khả năng giải quyết dựa trên ngữ cảnh(Solvability).} Tiêu chí này kiểm định khả năng câu hỏi có thể được giải quyết với văn bản đầu vào. Một câu hỏi đạt chuẩn tiêu chí này buộc phải được trả lời dựa vào nội dung của văn bản thay vì thông tin bên ngoài. Không chỉ vậy, câu hỏi phải có đúng một đáp án đúng duy nhất. Để tự động đánh giá độ đo này, nghiên cứu sử dụng LLM thương mại với chất lượng cao để trực tiếp giải quyết câu hỏi. Nếu kết quả của giám khảo giống như kết quả đã sinh, câu hỏi sẽ được đánh dấu đạt chuẩn và ngược lại. 
   $$\text{Solvability}(q) = \begin{cases} 1, & \text{nếu } C_{\text{judge}}(q) = \{c^*\} \\ 0, & \text{ngược lại} \end{cases}$$
    trong đó $C_{\text{judge}}(q)$ là tập hợp các phương án được LLM giám khảo chọn cho câu hỏi $q$, và $c^*$ là đáp án đúng duy nhất. Nếu chọn sai hoặc chọn nhiều phương án, câu hỏi bị đánh giá 0 (có thể do mơ hồ hoặc chứa nhiều đáp án đúng).

   \textbf{Độ bám sát phân loại nhận thức (Alignment):}  Tiêu chí này kiểm tra mức độ nhận thức hay độ khó của câu hỏi sinh ra có đúng với yêu cầu của hệ thống từ trước hay không. LLM giám khảo sẽ đối chiếu câu hỏi sinh ra với các tiêu chí của thang đo Bloom, từ đó đưa ra cấp độ thực tế của câu hỏi. Tương tự như tác tử \italic{Phân loại}, câu hỏi đạt tiêu chí này nếu cấp độ giám khảo dự đoán trùng với cấp độ đã tạo.
   $$\text{Alignment}(q) = \begin{cases} 1, & \text{nếu } l_{\text{judge}}(q) = l_{\text{target}}(q) \\ 0, & \text{ngược lại} \end{cases}$$
    trong đó $l_{\text{judge}}(q)$ là cấp độ do LLM dự đoán và $l_{\text{target}}(q)$ là cấp độ yêu cầu.

 \textbf{Chất lượng phương án nhiễu (Distractor Quality ):} Tiêu chí này đánh giá chất lượng của ba phương án nhiễu trong mỗi câu hỏi. Một phương án sai hợp lệ phải thỏa mãn đồng thời hai tiêu chí: trái ngược hoàn toàn so với đáp án đúng ((không tạo ra sự mơ hồ hay đa nghĩa cho người học) và có khả năng đánh lừa những học sinh không nắm vững kiến thức. LLM giám khảo sẽ đánh dấu các phương án nhiễu hợp lệ, giá trị của độ đo sẽ là tổng số phương án nhiễu hợp lệ trung bình đã được chuẩn hóa.

Gọi $v(q)$ là số phương án nhiễu hợp lệ của câu hỏi $q$ (với $v(q) \in \{0, 1, 2, 3\}$). Chỉ số chất lượng cho mỗi cấp độ $l$ được chuẩn hóa về khoảng $[0, 1]$:
$$\text{DQ}_l = \frac{1}{N_l} \sum_{q \in \mathcal{A}_l} \frac{v(q)}{3}$$
trong đó $\mathcal{A}_l$ là tập hợp các câu hỏi đã được chấp nhận ở cấp độ $l$, và $N_l = |\mathcal{A}_l|$.

Chỉ số Distractor Quality tổng thể là trung bình cộng trên ba cấp độ:
$$\text{DQ}_{\text{overall}} = \frac{1}{3} \sum_{l=1}^{3} \text{DQ}_l$$

Trong quá trình thực nghiệm, nghiên cứu phát hiện ra rằng một số mô hình yếu có xu hướng tạo ra các câu hỏi dễ (thường là cấp độ 1 và 2), dẫn đến việc kết quả {Alignment/] thấp trong khi \i{Solvability{ lại cao. Hay nói cách khác, các mô hình này sinh ra các câu hỏi có chất lượng tốt về hình thức (đúng một đáp án, không mơ hồ, đúng ngữ cảnh) nhưng không đúng cấp độ nhận thức yêu cầu. Để các đánh giá được công bằng và hiệu quả, nghiên cứu đề xuất độ đo Tỷ lệ chấp nhận (Acceptance Rate) là tỷ lệ các câu hỏi đồng thời thỏa mãn cả hai tiêu chí \italic{Solvability} và \italic{Alignment}. Đây cũng là độ đo chính được sử dụng trong báo cáo này.
$$\text{Acceptance}(q) = \text{Solvability}(q) \times \text{Alignment}(q)$$
$$\text{Acceptance Rate} = \frac{1}{N} \sum_{i=1}^{N} \text{Acceptance}(q_i)$$
trong đó $N$ là tổng số câu hỏi được đánh giá. Chỉ những câu hỏi có Acceptance bằng 1 mới được đưa vào Giai đoạn 2.

\subsection{Đánh giá hiệu năng nhận dạng văn bản}

Bài toán thứ hai cần được đánh giá của khóa luận là tác vụ OCR. Nghiên cứu sử dụng hai thang đo tiêu chuẩn là tỷ lệ lỗi ký tự (CER - Character Error Rate) và tỷ lệ lỗi từ (WER - Word Error Rate).Nguyên lý của hai độ đo này là dựa trên khoảng cách chỉnh sửa Levenshtein. Cụ thể, giá trị của độ đo là số các bước sửa đổi tối thiểu để biến đổi chuỗi dự đoán thành chuỗi tham chiếu gốc với ba thao tác chèn ($I$), xóa ($D$) và thay thế ($S$).

$$\text{CER} = \frac{S_c + D_c + I_c}{N_c}$$
$$\text{WER} = \frac{S_w + D_w + I_w}{N_w}$$

trong đó $N_c$ và $N_w$ lần lượt là tổng số ký tự và tổng số từ trong chuỗi tham chiếu. Giá trị CER và WER càng thấp thể hiện chất lượng nhận dạng OCR càng tốt.


